\chapter{Gingivitis}

\section*{Definition and Overview}
Gingivitis is a common and mild form of gum disease in which the gums (gingivae) become inflamed, red, swollen, and prone to bleeding. It is caused by the build-up of plaque, a sticky film of bacteria that forms constantly on the teeth. When plaque is not removed by regular brushing and flossing, the bacteria irritate the gums and trigger inflammation. Gingivitis is very common, and its hallmark is bleeding when brushing or flossing. It is reversible with good oral hygiene and professional cleaning, but if left untreated it can progress to periodontitis, a more serious infection that damages the tissues and bone supporting the teeth and can lead to tooth loss.

\section*{Epidemiology}
Gingivitis is one of the most widespread health conditions in the world, affecting the majority of adults at some point in their lives. In Australia, around half of adults have some form of gum disease, and gingivitis is particularly common in teenagers and young adults as well as during pregnancy, when hormonal changes increase gum sensitivity. Prevalence increases with age, and risk is higher in smokers, people with poorly controlled diabetes, and those with poor oral hygiene. Because gingivitis is usually painless, many people are unaware they have it until a dentist points it out.

\section*{Aetiology and Pathophysiology}
Gingivitis develops when bacteria in dental plaque trigger an inflammatory response in the gums.

\begin{itemize}
    \item \textbf{Plaque formation:} bacteria colonise the tooth surface within hours of brushing, forming a sticky film
    \item \textbf{Inflammation:} bacterial toxins stimulate the immune system, causing the gums to become red, swollen, and tender
    \item \textbf{Bleeding:} inflamed gums bleed easily because of increased blood flow and weakened tissue
    \item \textbf{Tartar (calculus):} if plaque is not removed, it hardens into tartar, which cannot be removed by brushing and provides a rough surface for more plaque
    \item \textbf{Contributing factors:} smoking, diabetes, stress, poor nutrition, hormonal changes, dry mouth, and certain medicines
    \item \textbf{Reversibility:} unlike periodontitis, gingivitis does not damage the bone and is fully reversible once plaque is removed
\end{itemize}

\section*{Clinical Features}
The signs of gingivitis are often noticed first by a dentist, but people may observe them themselves.

\begin{itemize}
    \item Gums that bleed during brushing or flossing
    \item Red, swollen, or puffy gums
    \item Gums that feel tender or sore
    \item Bad breath (halitosis) that persists despite brushing
    \item A bad taste in the mouth
    \item Receding gums, where the gums pull away from the teeth
    \item Gums that look shiny or bright red in patches
    \item Often no pain at all, which is why it can go unnoticed
\end{itemize}

\section*{Differential Diagnosis}
Other conditions can cause similar gum changes and should be considered, particularly when gingivitis does not respond to good hygiene.

\begin{itemize}
    \item Periodontitis, in which inflammation has spread to the bone supporting the teeth
    \item Necrotising ulcerative gingivitis, a painful, acute infection often associated with stress, smoking, and poor diet
    \item Viral infections such as herpes gingivostomatitis, which causes painful ulcers
    \item Fungal infection (oral thrush), which causes white patches
    \item Lichen planus and other inflammatory conditions of the mouth
    \item Side effects of medicines, including some blood pressure and epilepsy drugs that cause gum overgrowth
    \item Leukaemia and other blood disorders that can cause bleeding gums
    \item Vitamin C deficiency (scurvy), which causes bleeding gums
\end{itemize}

\section*{When to Seek Medical Care}
Everyone should have regular dental check-ups, at least every six to twelve months, to detect and treat gingivitis early. See a dentist sooner if gums bleed persistently, look swollen, or cause discomfort despite good brushing and flossing.

\textbf{Call 000 (or local emergency services) for:}
\begin{itemize}
    \item Sudden severe swelling of the face or jaw, especially with fever, which may indicate a dental abscess spreading
    \item Difficulty breathing or swallowing related to a dental infection
    \item Uncontrolled bleeding from the gums
    \item Any serious infection with fever and facial swelling
\end{itemize}

\section*{Diagnosis and Investigations}
Gingivitis is diagnosed by a dental examination, and most cases need no special tests.

\begin{itemize}
    \item \textbf{Clinical examination:} inspection of the gums for redness, swelling, and bleeding on gentle probing
    \item \textbf{Probing:} a periodontal probe measures pocket depth between the gum and tooth; depths of 3 mm or less with bleeding indicate gingivitis
    \item \textbf{Plaque disclosure:} a dye reveals areas where plaque has been missed
    \item \textbf{Dental X-rays:} not needed for gingivitis alone, but used to check for bone loss if periodontitis is suspected
    \item \textbf{Review of risk factors:} smoking, diabetes, medicines, and pregnancy status are considered
\end{itemize}

\section*{Management}
Gingivitis is managed by removing plaque and tartar and then maintaining excellent daily oral hygiene.

\begin{itemize}
    \item \textbf{Professional cleaning:} a dentist or oral health therapist removes plaque and tartar with scaling and polishing
    \item \textbf{Improved brushing:} brushing twice daily with fluoride toothpaste for two minutes
    \item \textbf{Flossing and interdental cleaning:} cleaning between the teeth daily with floss or interdental brushes
    \item \textbf{Antibacterial mouthwash:} chlorhexidine mouthwash may be recommended for short-term use
    \item \textbf{Follow-up:} reassessment to confirm the gums have returned to health
    \item \textbf{Addressing risk factors:} smoking cessation and good diabetes control improve outcomes
    \item \textbf{Regular maintenance:} routine professional cleaning at intervals suited to your gum health
\end{itemize}

\section*{Complications}
\begin{itemize}
    \item Progression to periodontitis, with irreversible damage to the supporting bone and eventual tooth loss
    \item Recurrent bleeding and bad breath that affect confidence and quality of life
    \item Gum recession and root exposure, which can cause tooth sensitivity
    \item Increased risk of infection if gums are damaged or procedures are performed on unhealthy gums
    \item Links to broader health: uncontrolled gum disease is associated with increased risk of cardiovascular disease, diabetes complications, and pregnancy problems
    \item Dental abscess in advanced disease
\end{itemize}

\section*{Prognosis and Outcomes}
The prognosis for gingivitis is excellent. With professional cleaning and consistent daily oral hygiene, inflamed gums usually return to normal within one to two weeks, and bleeding stops. Because gingivitis is fully reversible, the main challenge is sustaining the daily habits that prevent plaque from re-accumulating. People who maintain good oral hygiene after treatment keep their gums healthy and greatly reduce their risk of developing periodontitis and losing teeth later in life.

\section*{Prevention and Self-Care}
\begin{itemize}
    \item Brush your teeth twice a day for two minutes with fluoride toothpaste
    \item Clean between your teeth daily with floss or interdental brushes
    \item Use a soft-bristled toothbrush and replace it every three months or after illness
    \item Visit the dentist regularly, typically every six to twelve months
    \item Do not smoke, as smoking worsens gum disease
    \item Manage diabetes well, as high blood glucose increases gum disease risk
    \item Eat a balanced diet and limit sugary foods and drinks
    \item If your gums bleed when you start flossing, persist for a few days; if bleeding continues, see a dentist
\end{itemize}

\section*{Sources and Further Reading}
The following reputable sources provide further information for healthcare students and patients seeking reliable, current advice about gingivitis.

\begin{itemize}
    \item Australian Dental Association. \textit{Gum disease: information for patients.} https://www.ada.org.au/
    \item Healthdirect Australia. \textit{Gingivitis and gum disease.} https://www.healthdirect.gov.au/
    \item Dental Health Services Victoria. \textit{Teeth and gum care.} https://www.dhsv.org.au/
    \item NHS. \textit{Gum disease: symptoms and treatment.} https://www.nhs.uk/
    \item Mayo Clinic. \textit{Gingivitis: symptoms and causes.} https://www.mayoclinic.org/
    \item Australian Government Department of Health and Aged Care. \textit{Oral health and dental care.} https://www.health.gov.au/
\end{itemize}
